\chapter{Selecting events is not rewriting the history of events (in a continuous probability space)}
\section*{Abstract}
Using the simple fact described in the title, we prove the existence of
a computational problem with implications to Machine Learning, Quantum
Mechanics and Complexity Theory. We also prove P!=NP (the solution can
be verified in time polynomial in the number of bits of the input and
output (NP) but the problem cannot be solved in time polynomial in the
number of bits of the input and output (P)), but this claim still needs
to be reviewed by experts in Complexity Theory.


\hypertarget{introduction}{%
\section{1. Introduction}\label{introduction}}

In this article we will be using the words ``real'' as in
\(\mathbb{R}\), ``real-world'' and ``random'' to avoid misunderstandings
in contexts where we could (and perhaps should) just use the word
``real'' instead of ``real-world'' and the word ``non-deterministic''
(in the Physics sense, not in the Complexity Theory sense) instead of
``random''.

We can always select events with some feature without rewriting the
history of events, in a standard probability space. In fact, in a
continuous probability space we can select events such that a random
real variable \(y\in [0,1]\) (with probability given by the Lebesgue
measure) verifies \(y=0\), but there is no complete history of events
where \(y=0\) (always as would be required, or even just once) because
the probability space is continuous by assumption and thus the event
\(y=0\) has null probability (only an interval would have non-null
probability).

In this article we will show that this simple but non-trivial fact has
profound implications not only to Complexity
Theory\cite{novak2008tractability}\cite{hennig2022probabilistic}\cite{ritter2000average}\cite{bna} but also to Machine
Learning\cite{hennig2022probabilistic} and Quantum Mechanics (independently of the
implications to the P vs. NP problem). Almost all (in the sense we will
define in this article) real functions of a real variable cannot be
computed for all practical purposes, not even approximately. But a
random selection allows computations in polynomial-time complexity
involving the incomplete knowledge about a real function that cannot be
computed in polynomial-time complexity. This is a fundamental reason why
we cannot exclude a random time-evolution: a deterministic
time-evolution may exist, but it has so much complexity that it cannot
be calculated for all practical purposes, not even approximately (since
\(L^\infty\) is non-separable).

Note that whenever we deal with a non-separable space, there are issues
with computability\cite{computable} because some elements of a non-separable space cannot be
approximated by a finite set of elements, up to an arbitrarily small
error. For instance, \(L^\infty([0,1])\) and its dual space are both
non-separable\cite{book}. While there
are separable spaces of real functions of real variables, whenever we
add uncertainties/probabilities to such spaces (which is often required
when doing approximations) we tend to create non-separable
spaces\cite{book}, unless equivalence
relations change the space of functions. For instance, the set \([0,1]\)
is separable but the Lebesgue measure imposes that the rational numbers
in \([0,1]\) can be discarded, despite that the sets \([0,1]\) including
/excluding rational numbers are different. In the same way, the smooth
functions (or functions computable in a reasonable time) may be
discarded from a space of functions for particular
uncertainties/probabilities.

Using the simple fact mentioned above, we will also prove the existence
of a computational problem (defined by a continuous probability space)
whose solution can be verified in time polynomial in the numbers of bits
of the input and output (NP) but cannot be solved in time polynomial in
the numbers of bits of the input and output (P), when using only a
deterministic Turing machine\cite{pvsnp}\cite{10.1145ux2f800152.804913}. That is, P\(\neq\)NP.

The goal of this article is to define the specific problem
unambiguously, and the level of mathematical details will be adjusted to
that: too much mathematical detail would shift the focus from the
specific problem. Less detail does not always imply less mathematical
rigor. Since the present author is an expert in Physics but not an
expert in Complexity Theory, we will also try to prove the P vs. NP as
much as it is possible, but only as a secondary goal knowing that much
work by experts in complexity theory is still required because it is
likely that: 1) something went wrong in the relation described here
between the specific problem and the P vs. NP problem; and/or 2) the
specific problem indeed can be used to solve the P vs. NP problem, but
the proof presented here is incomplete.

\hypertarget{complexity-theory-in-the-context-of-probability-theory}{%
\section{2. Complexity Theory in the context of probability
theory}\label{complexity-theory-in-the-context-of-probability-theory}}

Complexity Theory can be studied in the context of probability
theory\cite{novak2008tractability}\cite{hennig2022probabilistic}\cite{ritter2000average}\cite{bna}, because many
real-world problems require approximations and uncertainties not only
due to the limitations of any computer (already accounted for by
Complexity Theory when using only finite numbers of bits, although there
is also room for improvement here) but also due to the limitations of
the measuring devices of physical phenomena and limitations of the
mathematical models used to approximate the real-world, for instance
when dealing with real variables. Most uncertainties are not related to
the computer used and do not get smaller when increasing the number of
bits in the computer. In fact, Physics as a science tries to be
independent of Computer Science and vice-versa, as much as possible.
Probability theory is a language (or interface) that allows us to
transfer a problem between two sciences (these two or others).

There are two possible approaches to errors or
uncertainties\cite{novak2008tractability}\cite{hennig2022probabilistic}\cite{ritter2000average}\cite{bna}: average error (with
respect to a probability measure) and maximal (except in sets of null
measure with respect to a probability measure) error. It turns out that
both approaches can be defined using Hilbert spaces: the average error
(that is, \(L^2\) norm) is defined by a normalized wave-function (an
element of the Hilbert space) being the square-root of a probability
density function; and the maximal error (that is, \(L^\infty\) norm) is
defined by an element of the abelian von Neumann algebra of operators on
the Hilbert space.

The average error is relevant because even if P!=NP it could still make
no difference with respect to the scenario P=NP for many practical
purposes, if every NP problem with a reasonable probability distribution
on its inputs could be solved in polynomial time-complexity on average
on a deterministic Turing
machine\cite{avcomplete}\cite{impagliazzo1995personal}.

Moreover, since some functions are real constant functions, the
definition of a real function must be consistent with the definition of
a real number. Certainly, a natural definition of a real number in the
context of Complexity Theory uses a standard probability space.
Non-standard probability spaces are rarely (or never) used in
Experimental Physics, so it is not obvious how useful a ``real number''
(or ``real function'') defined in a non-standard probability space could
be in real-world applications. There are only countable or continuous
measures (or mixed) in a standard probability space, then we can define
exactly only a countable number of real numbers (usually the rationals,
but not necessarily), the remaining real numbers can only be constrained
to be inside an interval with a finite width, eventually very small but
never zero. We should also define all real functions using a standard
probability space, unless we find a fundamental reason not to do it (we
will not find it in this article).

We require a continuum standard probability space, since we can always
define a regular conditional probability density which implements a
selection of events in such probability
space\cite{durrett2019probability}. For
instance, this is what we do when we neglect the intrinsic computation
error (due to cosmic rays and many other reasons), we define a
deterministic function by selecting only certain events from a complete
history of random events. It is well known since many decades that in an
infinite-dimensional sphere of radius \(1\) (subset of a real Hilbert
space) there is a uniform prior measure induced by the \(L^2\)-distance
in the Hilbert
space\cite{Peterson2019GaussianLA}.
Every point in the sphere has null measure, only regions of the sphere
with non-null distance between some of its points may have non-null
measure (compatible with the uniform prior measure). This implies that
any knowledge (compatible with the uniform prior measure on the sphere)
about a real normalized wave-function has necessarily uncertainties,
defined by a connected region in the sphere with non-null maximum
\(L^2\) distance to some wave-function, however small it might be.

The discrete nature of the Turing machine is certainly compatible with a
continuous probability space: the number of bits of the input or output
can be arbitrarily large, and it is proportional to the logarithm of the
resolution of the partition of the interval \([0,1]\), with each
disjoint set of the partition corresponding to a different binary
number. Excluding a continuous cumulative prior would be unjustified,
for many reasons including: no prior is better for all
cases\cite{dutchbook}, there
are many problems where a step cumulative prior would not fit well (for
instance there is no uniform measure for the rationals in \([0,1]\) only
for the reals); it is hard to formulate any real-world problem where
only a step cumulative prior is used (think about the numbers \(\pi\) or
\(\sqrt{2}\) in numerical approximations, for instance), we usually use
a mixture of step and continuous cumulative priors; we can map an
ensemble of discrete random variables one-to-one to the real numbers,
for instance an ensemble of fair coins corresponds to the uniform real
measure in the interval \([0,1]\); also any real-world computer has an
intrinsic computation error (due to cosmic rays and many other reasons)
which is usually very small, but it cannot be eliminated. Thus, while we
can formulate a new unsolved version of the P vs. NP problem where only
a step cumulative prior is accepted, such version of the problem has
little to do with real-world computers and real-world problems.

Note that there are \(2^{2^n}\) different boolean\footnote{That is, a
  function with {{\(m = 1\)}{{{}{m}{}{=}{}}{{}{1}}}} boolean outputs}
functions on \(n\) boolean variables, Shannon proved that almost all
Boolean functions on \(n\) variables require circuits of size
\(\mathcal{O}(2^n/n)\)
\cite{shannon1949synthesis}\footnote{See
  also:
  \url{https://math.stackexchange.com/questions/756813/do-there-exist-polynomials-not-computable-in-polynomial-time}},
thus the time complexity of almost all Boolean functions on \emph{n}
variables for a Turing machine is at least (and at most, for all Boolean
functions)
\(\mathcal{O}(2^n/(n\log(n)))\)\cite{10.1145ux2f322123.322138} which is not polynomial in \(n\). Thus, almost all
numerical functions are not in the complexity class \(P\), according to
the uniform prior measure for any resolution of the partition. Moreover,
the uniform prior measure is compatible with a prior measure which
excludes (both in the maximal and in the average error approaches) all
real functions which are approximated by numerical functions with
complexity class \(P\), for all resolutions bigger than some resolution
of the partition.

\hypertarget{definition-of-the-problem}{%
\section{3. Definition of the problem}\label{definition-of-the-problem}}

Consider three measure spaces \(X=Y=[0,1]\in \mathbb{R}\) and
\(X\otimes X\) with the Lebesgue measure, corresponding to inputs (\(X\)
or \(X\otimes X\)) and an output (\(Y\)) of real functions. Given an
input in \(X\), we define a regular conditional probability density
which is a function from \(X\otimes X\to Y\), given by the probability
that a function for an input in \(X\) some constant output \(y\in Y\).
But there is always also a marginal probability density for \(Y\), and
we cannot say without uncertainty which is the output, because the
corresponding prior probability density would be incompatible with the
prior Lebesgue measure (by the Radon--Nikodym theorem). Thus, the
regular conditional probability density is a deterministic selection of
events which cannot be a complete history of events.

In a standard measure space it is always possible to define regular
conditional
probabilities\cite{durrett2019probability} and to choose the probability density \(p(x)=p_0(x)>0\) for all
\(x\in X\), except in sets with null measure. Thus, we will define
\(p(x\otimes y)=p(y|x)p_0(x)\) the joint probability density for the
tensor product \(X\otimes Y\) for a particular \(p(x)=p_0(x)>0\) for all
\(x\in X\), except in sets with null measure. Then, we can obtain any
other joint probability density \(p(x\otimes y)\) from the expression
\(p(y|x)p(x)=\frac{p(x\otimes y)}{p_0(x)}p(x)\).

The following results are valid for a random input in the interval
\([0,1]\) (which is a standard probability space) and also for an input
(or output) without uncertainties up to sets with null measure with
respect to the prior marginal measure of the input (or output), because
we use regular conditional probabilities (which always exist in standard
probability
spaces\cite{durrett2019probability})
for fully known inputs (or outputs). This is crucial, since the input
includes two samples from a uniform distribution in \([0,1]\) which may
generate numerical functions in \(P\) when the sample is in a set of
null measure.

However, a (continuous cumulative) probability distribution does not
contain enough information to unambiguously define a function. On the
other hand, a real wave-function whose square is the joint probability
distribution allows the definition of a unitary operator on a separable
Hilbert space. A unitary operator is a random generalization of a
deterministic symmetry transformation of a (countable or continuous)
sample space. Any unitary operator defined by a wave-function of two
continuous variables cannot be a deterministic symmetry transformation
(for similar reasons that a continuous probability distribution cannot
unambiguously define a function).

Since \(p(x\otimes y)\geq 0\) then there is always a normalized
wave-function \(\Psi\in L^2(X\otimes Y)\) such that
\(|\Psi(x\otimes y)|^2=p(x\otimes y)\). Note that the Koopman-von
Neumann version of classical statistical
mechanics\cite{Sudarshan1976Interaction} defines classical statistical mechanics as a particular case of
quantum mechanics where the algebra of observable operators is
necessarily commutative (because the time-evolution is deterministic).
In an infinite-dimensional sphere of radius \(1\) (subset of a real
Hilbert space) there is a uniform prior measure induced by the
\(L^2\)-distance in the Hilbert space. We choose a prior measure
(compatible with the uniform prior measure) which excludes all real
functions which are approximated by numerical functions with complexity
class \(P\), for a high enough resolution of the partition.

Given an input in \(([0,1])^2\) (the input consists of two samples from
a uniform distribution in the interval \([0,1]\), imported from ANU QRNG
\cite{PhysRevApplied.3.054004}\footnote{See also: \url{https://qrng.anu.edu.au/random-binary}}
for instance) and a candidate output in \([0,1]\) the wave-function
uniquely defines a random symmetry transformation. Such symmetry
transformation is not lacking information since it can be inverted (the
non-unitary isometries have null prior measure). The cumulative
probability distribution is given by the integral of the modulus squared
of the wave-function in the corresponding region of the sample space.

From the cumulative marginal probability distribution, such that \(Y\)
is fully integrated we determine \(x\in X\) using the first sample from
the uniform distribution in \([0,1]\). From the cumulative conditional
probability distribution with the condition that \(x\in X\) is what we
determined previously, we determine \(y\in X\) using the second sample
from the uniform distribution in \([0,1]\). We apply the
inverse-transform sampling
method\cite{Rubinstein2016Simulation}, that is check in the interval of the partition defined by
the bits corresponding to \(X\) and \(Y\), whether the cumulative
distribution crosses the sample from the uniform distribution. This
defines the deterministic verification of the candidate output
corresponding to the input, in agreement with the Born's rule of Quantum
Mechanics.

Note that the resulting deterministic function is not necessarily
invertible, because the collapse of the wave-function is irreversible
(unless the symmetry transformation would be deterministic, which is
excluded in this case because we have a continuous probability space of
functions).

\hypertarget{the-classical-turing-machine-defined-as-a-quantum-computer}{%
\section{4. The classical Turing machine defined as a Quantum
computer}\label{the-classical-turing-machine-defined-as-a-quantum-computer}}

The Turing machine can be equivalently defined as the set of general
recursive functions, which are partial functions from non-negative
integers to non-negative
integers\cite{sep-recursive-functions}.
But the set of all functions from non-negative integers to non-negative
integers is not suitable to define a measure, since they form an
uncountable set, in a context where the continuum is not defined.
Moreover, the general recursive functions are based in the notion of
computability (that the Turing machine halts in a finite time), but
computability does not hold in the limit of an infinite number of input
bits, thus to study such limit we need to define uncomputable functions
somehow (we will use complete spaces, where Cauchy sequences always
converge to an element inside the space).

On the other hand, it is widely believed (and we will show in the
following) that any computational problem that can be solved by a
classical computer can also be solved by a quantum computer and
vice-versa. That is, quantum computers obey the Church--Turing thesis.
Note that it is well known that some circuits (classical hardware)
provide exponential speedups when compared with some other circuits in
some functions (because the input bits can be reparametrized, this is
why the time complexity of a function has an upper bound, but it is not
known how to establish a lower bound; it is also consistent with the
fact that the halting problem is undecidable, that is, given an
arbitrary function from integers to integers and an arbitrary input, we
cannot determine if the output of such function is computable or not),
thus the fact that a classical Turing machine can be defined as a
Quantum computer is compatible with the fact that quantum computers
provide exponential speedups when compared with some classical computers
in some functions.

We start by noticing that the domain of a general recursive function can
be defined by a dense countable basis of a particular Hilbert space
which is the (Guichardet) \(L^2\) completion of the set of all finite
linear combinations of simple tensor products of elements of a countable
basis of a base Hilbert space, where all but finitely many of the
factors equal the vacuum
state\cite{bratteli1987operator} (like in a Fock-space, but without the symmetrization). But the
unitary linear transformations on a normalized wave-function are not
necessarily the most general transformations of the probability measure
corresponding to the wave-function. Because of that, we build a
Fock-space where the base Hilbert space is the previous
Guichardet-space, then the unitary transformations on this
Fock-Guichardet-space allow us to implement the most general
transformations of a probability measure, corresponding to a normalized
wave-function in the base Guichardet-space. Note that a countable basis
of the Guichardet-space is already made of simple tensor products, and
the simple tensor product is associative, thus the Fock-Guichardet-space
is isomorphic to the Guichardet-space, but we still prefer to use the
Fock-Guichardet-space due to the existence of standard tools for
Fock-spaces.

Note that the input of a general recursive function is a finite number
of integers, but its output is only one integer. However, any function
which outputs several integers is a direct sum of functions which output
one integer. The other way around is also true, once we define a vacuum
state (included in the Fock-Guichardet-space), that is, a function which
outputs one integer is a particular case of a function that outputs
several integers where all outputs except one correspond to the vacuum
state. Thus, we can consider only unitary automorphisms of the
Fock-Guichardet-space.

To be able to define a measure, we make the integers correspond to the
(countable) step functions with rational endpoints in the interval
\([0,1]\) and weights which are plus or minus the square root of a
rational number\cite{book91891579}. The vacuum state is the constant positive function
with norm 1, and it corresponds to the integer \(0\). We eliminate
duplicated step functions in the correspondence with the integers, for
instance if two neighbor intervals have the same weight then they are
fused.

Then, the limit of infinitesimal intervals is well-defined, and it is
defined by an element of \(L^2([0,1])\). Since the general recursive
functions are partial functions, then they are a particular case of
partially-defined linear operators \(L^2([0,1])\to L^2([0,1])\), and we
can define the base Hilbert space of the Fock-Guichardet-space as
\(L^2([0,1])\).

\hypertarget{worst-case-prior-measure-rational-functions-and-radical-determinism}{%
\section{5. Worst case prior measure, rational functions and radical
determinism}\label{worst-case-prior-measure-rational-functions-and-radical-determinism}}

In a standard probability space, there are only continuous and/or
countable measures. However, these may be mixed in an arbitrary way. For
a theory of Physics we could choose the best case prior measure (as we
did in the previous sections), since we just want to find a prior which
is consistent with the experimental data, without many concerns about
alternative priors. However, in Cryptography we need guarantees that our
limits are robust under arbitrary choices, so we need to assume the
worst case prior measure.

The previous sections could also be made compatible with the worst case
prior measure, if we had a computer capable of comparing real numbers
not rational. That would be acceptable for a theory of Physics, but it
would make it difficult to obtain guarantees for Cryptography.

It is also difficult to guarantee true randomness is real-world
applications of Cryptography. Since probabilities only mean incomplete
information, we can use Probability Theory in the context of radical
determinism (where nothing is random). For Cryptography, we need the
worst case prior measure, rational functions and radical determinism.

So we start by eliminating a non-standard probability measure: any
probability theory is a universal language (like English or mathematical
logic) to define abstract models of the objects we want to study. A
standard probability theory is universal and irreducible, meaning that
it has the minimal content to be considered a probability theory (in
agreement with Quantum Mechanics and Experimental Physics, for
instance). If the non-standard probability theory is also irreducible,
then the corresponding models are equivalent, and we can use the
standard version without loss of generality. This allows us to transfer
models between different sciences. But often the non-standard
probability theory is reducible, this means that the boundary between
model and probability theory is not where it would be in the standard
case and there are properties that we are attributing to the probability
theory that in fact belong to the model.

Thus, we should assume a standard probability theory and leave some
flexibility in the definition of the computer model and not the other
way around, as it often happens in Complexity Theory where there are
strict axioms for different computer models, while asymptotic limits are
taken without defining the probability space, which is recipe to end up
with mathematical results and questions which are hard to transfer to
experimental physics and many other sciences.

In the context of radical determinism, the history of events is a
non-random countable sequence of events. Thus, some events with null
measure might happen, due to radical determinism. But only a countable
number of such events. That means that a continuous measure is only
truly continuous up to a countable number of points, this possibility is
already considered by the worst case prior measure.

Consider now a boolean function of a countable infinite number of bits.
The time complexity of almost all Boolean functions on \emph{n}
variables for a Turing machine is at least (and at most, for all Boolean
functions)
\(\mathcal{O}(2^n/(n\log(n)))\)\cite{shannon1949synthesis} \cite{10.1145ux2f322123.322138} which is not polynomial in \(n\), but the same boolean
function has different time complexities in different
circuits\cite{shannon1949synthesis}
(because the input bits can be reparametrized). Thus, the time
complexity of a function depends on the circuit (computer design). Also,
an algorithm with polynomial time-complexity is not guaranteed to be
fast (due to large order and coefficients of the polynomial), thus only
the asymptotic behavior is fast, and so we cannot put an upper bound on
the number of bits of the input. But the arbitrarily large number of
bits of the input introduces another ambiguity: given any problem with
any time complexity (exponential \(\mathcal{O}(2^n)\), for instance)
there is a problem with linear time-complexity in the number of bits
that takes the same amount of time to run. Thus, the polynomial
time-complexity is faster than exponential time-complexity in asymptotic
behavior, only for the same number of bits of the input. This is a
condition without much meaning when the input has a countable infinite
number of bits.

What we can say is that the countable (or mixed) measure allows defining
functions that eventually have polynomial time complexity (that is, not
necessarily non-polynomial time complexity). This contrasts with the
necessarily non-polynomial time complexity of the functions defined by
the continuous measure. In the following we will show that, given any
mixed prior measure, we can always redefine the problem to have a
continuous prior measure such that its results cannot be reproduced by a
mixed prior measure, effectively converting the worst case prior measure
into the best case prior measure.

\hypertarget{converting-the-worst-case-measure-to-best-case-measure}{%
\subsection{\texorpdfstring{\href{}{\emph{Converting the worst case
measure to best case
measure}}}{Converting the worst case measure to best case measure}}\label{converting-the-worst-case-measure-to-best-case-measure}}

The prior measure must be mixed or continuous, to allow for the limit of
an infinite number of bits of all computable functions. But we must
prove that \(P\neq NP\) in a single deterministic Turing machine (one
with a continuous prior measure, for instance), not in the set of all
possible deterministic Turing machines. That would not be possible,
since for any computable function \(f\) (including any function in
\(NP\)) there is a deterministic Turing machine where \(f\) is in \(P\),
by reparametrizing the input bits.

It is not obvious if we can prove \(P\neq NP\) in any single
deterministic Turing machine, or just in one particular deterministic
Turing machine. However, given a worst case prior measure (thus mixed
measure), there is a subset of the input random sample which also
implements a deterministic Turing machine and where the prior measure is
continuous, where \(P\neq NP\), thus it becomes legitimate to claim that
this fact already shows that \(P\neq NP\). Moreover, using subsets of
the input random sample (thus, regular conditioned probability) we can
create any other prior measure, because any abelian von Neumann algebra
of operators on a separable Hilbert space is *-isomorphic to exactly one
of the following:

\begin{itemize}
\item
  \(l^\infty(\{ 1 , 2 , ... , n \} ) , n \geq 1\)
\item
  \(l^\infty(\mathbb{N})\)
\item
  \(L^\infty([0,1])\)
\item
  \(L^\infty([0,1]\cup\{ 1 , 2 , ... , n \} ) , n \geq 1\)
\item
  \(L^\infty([0,1]\cup \mathbb{N})\).
\end{itemize}

Equivalently, a standard probability space is isomorphic (up to sets
with null measure) to the interval \([0,1]\) with Lebesgue measure, a
finite or countable set of atoms, or a combination (disjoint union) of
both.

Thus, for the worst case prior measure if we include all subsets of the
input random sample, then using two integers (which is countable) we
include a countable set of deterministic Turing machines \(\{k\}\) and
countable functions \(f_{k,n}\), one machine for each countable function
\(f_{k,1}\) such that \(f_{k,1}\) is in \(P\) and \(f_{k,n}=f_{n,k}\),
then we cannot prove \(P\neq NP\) (in fact we would conclude that
\(P=NP\), due to the possibility of reparametrizing the input bits).

Then, we can only prove \(P\neq NP\) in one particular deterministic
Turing machine and not in any single deterministic Turing machine.

Given any mixed prior measure, there is an interval of the input random
sample where it is continuous. We rescale to \([0,1]\) all intervals
corresponding to such interval where the mixed measure is continuous,
using conditioned probability. The results in (the new interval of the
input random sample) \([0,1]\) cannot be fully reproduced by any other
measure which is mixed in \([0,1]\).

Given any other measure which is mixed in \([0,1]\), there is an
interval with rational endpoints of the input random sample where there
is a finite difference between the two cumulative probability
distributions, otherwise both would be continuous. This translates into
two different averages of \(y\) in an interval (for \(x\)) of a
partition of \([0,1]\), separated by a finite difference.

Then the indicator function of a \(y\) corresponding to the continuous
measure, in the interval (for \(x\)) of the partition of \([0,1]\), is
in \(NP\) (more precisely, it can be extended to be in \(NP\), since we
only defined for \(x\) in a subset of \([0,1]\)) but it cannot be
reproduced by the mixed measure. It cannot be reproduced by the
continuous prior measure either, since a function constant in \(x\) in a
finite interval has null measure (for a continuous prior measure). Note
that the function that corresponds to the indicator function above
defined, is a function constant in \(x\) in a finite interval.

Note that a continuous prior measure admits a regular conditional
probability density, which allows us to define a selection
(verification) of a candidate output. The verification of a constant
output is in \(NP\) and thus in a mixed measure, but it requires one
more input (the candidate output) and thus it is compatible with a
continuous measure of functions of \(x\). That is, the measure is
overall mixed, being continuous only for functions with one input.

\hypertarget{pneqnp}{%
\section{\texorpdfstring{6.
P\(\neq\)NP}{6. P\textbackslash neqNP}}\label{pneqnp}}

When using a computer to solve the problem defined in the previous
section, any disjoint set of the partition is an interval. This gives us
two options and two options only, either the selection of events is
fully deterministic or only approximately deterministic:

\begin{enumerate}
\item
  The selection of events is fully deterministic. Then we impose a
  condition on the output \(Y\) of any wave-function in the sphere, this
  defines a regular conditional probability density for the input \(X\)
  conditioned on a constant rational output \(y\). As shown in the
  previous section, there is a rational \(y\) corresponding to the
  continuous measure, in an interval with rational endpoints (for \(x\))
  of the partition of \([0,1]\), which it cannot be reproduced by a
  countable prior measure. It cannot be reproduced by the continuous
  prior measure either, since a function constant in \(x\) in a finite
  interval has null measure (for a continuous prior measure). That is,
  there is no function in \(P\) corresponding to the indicator function
  for \(y\), which is in \(NP\) (more precisely, it can be extended to
  be in \(NP\), since we only defined it for \(x\) in a subset of
  \([0,1]\)). This implies \(P\neq NP\).
\item
  The selection of events has some randomness, as small as we want. Then
  we do not impose a condition on the output (eventually we impose a
  condition on the input \(X\), depending on whether we want fixed or
  averaged input). In a strict interpretation of the \(P\) vs. \(NP\)
  problem this option is excluded by definition since the official
  formulation assumes that both the verification and the solution of the
  problem are both fully deterministic. This already implies
  \(P\neq NP\), in a strict interpretation.
\end{enumerate}

Note that a complete history of events needs to be countable, so that we
can convert it into a single event (mapping complete histories of events
one-to-one to the real numbers in the interval \([0,1]\), for instance).
We could also define a density (that is, yet to be integrated, using the
disintegration
theorem\cite{chang1997conditioning}) of an event. Such density is a regular conditional probability,
since regular conditional probabilities always exist in standard
probability
spaces\cite{durrett2019probability}.
But a density cannot correspond to a single event (by definition) and
thus it cannot be considered a complete history of events.

This proof is dependent on the fact that the prior measure is
continuous. If it in part continuous, and in part countable, then we can
choose just the continuous part of the sample space (see the previous
section). While we can use a countable part of the sample space to
approximately solve a continuous problem, and a continuous part of the
sample space to solve a countable problem, we cannot change the prior
measure from continuous to countable or vice-versa (by the Radon-Nikodym
theorem), because there is no Radon-Nikodym derivative between the two
measures, since the sets of null measure are disjoint between the two
measures. The prior measure defines the physical world where the
computer exists, thus it cannot be removed from any complete computer
model related to a physical computer.

Note that in the first paragraph of the official statement of the P vs.
NP problem\cite{pvsnp}, it is stated:

\begin{quote}
\emph{To define the problem precisely it is necessary to give a formal
model of a computer. The standard computer model in computability theory
is the Turing machine, introduced by Alan Turing in 1936.}
\emph{Although the model was introduced before physical computers were
built, it nevertheless continues to be accepted as the proper computer
model for the purpose of defining the notion of computable function.}
\end{quote}

As for any other proof, this proof is only as good as the axioms used
(that is, assumptions). The computer model used for a solid proof of the
P vs. NP problem should be widely accepted as a good approximation to a
physical computer for the purpose of defining the notion of computable
function. We believe our computer model is accepted by most experts in
Physics (as argued in the previous section). We claim that our computer
model makes no more assumptions than those required by the official
statement\cite{pvsnp} (including the
deterministic Turing machine), and it is as close to a physical computer
as possible, by today standards. Assuming a countable prior measure is
not realistic (as argued in the previous sections, for instance it would
exclude an ensemble of fair coins).

However, we believe that allowing a random selection of events is even
more realistic (as discussed in the previous sections, also with
implications to Machine Learning and Quantum Mechanics). In the next
section, we will define a selection of events which has some randomness
(as small as we want) and prove that even in that case, we still have
P\(\neq\)NP.

\hypertarget{realistic-version-of-the-problem-still-pneqnp}{%
\section{\texorpdfstring{7. Realistic version of the problem (still
P\(\neq\)NP)}{7. Realistic version of the problem (still P\textbackslash neqNP)}}\label{realistic-version-of-the-problem-still-pneqnp}}

A selection of events which is only approximately deterministic can be
approximated by a step function (step functions are dense in \(L^2\))
and thus there is a square with non-null constant measure. We rescale
such square to \([0,1]\times[0,1]\). We then consider a real polynomial
wave-function that is near the point in the sphere corresponding to a
constant wave-function, up to an error in the \(L^2\) norm which can be
as small as we want because the polynomials are dense in \(L^2\) (the
corresponding numerical polynomial does not need to be in \(P\)).

The first sample from the uniform distribution defines directly
\(x\in X\). An approximation (in the \(L^2\) norm) with polynomial
time-complexity to the selection function, is defined by setting
\(y\in Y\) equal to the second sample from the uniform distribution.

Since the wave-function is polynomial non-constant, then the
corresponding cumulative probability distribution minus the second
sample is strictly crescent (except in sets of null measure). Thus, when
we define the corresponding deterministic function we can choose the
second sample which produces an output which is as far from zero as we
want (in the interval \([0,1]\)), because we are using the \(L^\infty\)
norm now. We cannot average over the random sample, otherwise we need a
random computer (see next section). Thus, no approximation is possible,
and it suffices that we define a partition for the output which has two
disjoint sets (the measures of the sets are arbitrary, as long as they
are non-null) and a numerical output with one bit. Then, almost all
numerical functions are not in the \(P\) class, according to the prior
measure.

\hypertarget{generation-of-random-numbers-has-linear-time-complexity}{%
\section{8. Generation of random numbers has linear
time-complexity}\label{generation-of-random-numbers-has-linear-time-complexity}}

The two random samples from a uniform distribution in the interval
\([0,1]\) are inputs, in the deterministic Turing machine. However, in
the real-world these samples need to be generated somewhere and in
polynomial time-complexity, otherwise the time complexity of the random
selection computed by the real-world random computer could be
non-polynomial in the number of bits of the random samples.

Moreover, it would be better if the generation of random samples had
linear time complexity, since then we could do a constant rate of
experiments over time to validate the probability distribution of the
random selection, otherwise it would be impractical to generate an
infinite sequence of experiments.

We cannot prove this mathematically (since we would need more axioms).
However, the implications of this article to Quantum Mechanics help to
clarify the source of randomness of Quantum Mechanics (and thus of the
random samples). It is relevant to verify empirically that the
generation of random numbers with linear time complexity is possible,
for all practical purposes. We can visually check on the website from
ANU QRNG \footnote{See also: \url{https://qrng.anu.edu.au/random-binary}}
that the number of bits of the random sample grows linearly in time and
any complete history of events converges to a uniform probability
distribution.

Moreover, the entropy is maximal, in the sense that the deterministic
function needed to correlate the bits is not computable for all
practical purposes, not even approximately (since \(L^\infty\) is
non-separable) according to the prior measure.

\hypertarget{on-the-consequences-to-machine-learning}{%
\section{9. On the consequences to Machine
Learning}\label{on-the-consequences-to-machine-learning}}

In the introduction we discussed the implications (of the results of
this article), which are common to Machine Learning and Quantum
Mechanics. But Machine Learning (for instance Deep Neural Networks) is
not firmly based in probability theory, unlike Quantum Mechanics, then
there are more consequences.

In Machine Learning, methods inspired by probability theory are used
often\cite{hennig2022probabilistic}, but the formalism is based in approximations to
deterministic functions, guided by a distance (or equivalently, an
optimization problem) and not a measure. In fact, two of the main open
problems are the alignment of models and the incorporation of prior
knowledge\cite{muller2022transformers}, which could be both well solved by a prior measure if there
would be any measure defined.

Our results imply that under reasonable assumptions, almost all
functions are not computable not even approximately. Thus, Machine
Learning works because the functions we are approximating are in fact
probability distributions (eventually after some
reparametrization\cite{anonymous2023value-probability}). This shouldn't be surprising, since Classical Information
Theory shows (under reasonable assumptions) that probability is
unavoidable when we are dealing with representations of
knowledge/information\cite{novak2008tractability}\cite{hennig2022probabilistic}. But in Machine Learning the probability measure is
not consistently defined (despite that many methods are inspired by
probability theory), the probability measure emerges from the
approximation\cite{anonymous2023value-probability} and often in an inconsistent way. The inconsistency is not due to
a lack of computational power since modern neural networks can fit very
complex deterministic functions and fail
badly\cite{lakshminarayanan2017simple}\cite{degrancey:hal-03769683} in relatively simple probability distributions (e.g. catastrophic
forgetting or the need of calibration to have some probabilistic
guarantees\cite{degrancey:hal-03769683}).

This unavoidable emergence of a probability measure should be
investigated as a potential source of inefficiency, inconsistency and
even danger. If the emergence of a probability measure is unavoidable,
why don't we just define a probability measure in the formalism
consistently? Many people say ``it is how our brain works'', so
mathematics should step aside when there is empirical evidence.

But the empirical evidence is: oversized deep neural networks still
generalize well, apparently because often the learning process converges
to a local maximum (of the optimization problem) near the point where
the learning begun\cite{li2018learning}. This implies that if we repeat the learning process with a
random initialization (as we do when we consider ensembles of neural
networks\cite{lakshminarayanan2017simple}\cite{egele2022autodeuq}),
then we do not expect the new parameters to be near any particular
value, regardless of the result of the first learning process. This
expectation is justified by the fact that every three layers of a wide
enough neural network is a universal approximator of a
function\cite{hornik1989multilayer}, so any deviation introduced by three
layers can be fully corrected in the next three layers, when composing
dozens or hundreds of layers as we do in a deep neural network. Then the
correlation between the parameters corresponding to different local
maximums converges to zero, when the number of layers increases.

Thus, there is empirical evidence that oversized deep neural networks
still generalize well, precisely because a prior measure emerges: deep
learning does not converge to the global maximum and instead to one of
the local maximums chosen randomly, effectively sampling from a prior
measure in the sample space defined by all local maximums. This is
consistent with the good results achieved by ensembles of neural
networks\cite{lakshminarayanan2017simple}\cite{egele2022autodeuq},
which mimic many samples. However, it is a prior measure which we cannot
easily modify or even understand, because the measure space is the set
of all local maximums of the optimization problem. But, since we expect
the parameters to be fully uncorrelated between different local
maximums, then many other prior measures (which we can modify and
understand, such as the uniform measure) should achieve the same level
of generalization.

This is not a surprise, since oversized statistical models that still
generalize well were already found many decades ago by many
people\cite{Peterson2019GaussianLA}: a
standard probability space with a uniform probability measure can be
infinite-dimensional (the
sphere\cite{Peterson2019GaussianLA}
studied in this article, for instance).

More empirical evidence: no one looks to a blurred photo of a gorilla
and says with certainty that it is not a man in a gorilla suit. We all
have many doubts, when we are not sure about a subject we usually
express doubts through the absence of an action (not just us, but also
many animals), for instance we don't write a book about the subject we
don't know about.

There is no empirical evidence that our brain tries to create content
which is a short distance from content (books, conversations, etc.)
created under exceptional circumstances (when doubts are minimal). When
we are driving, and we do not know what is in front of us, we usually
just slow down or stop the car. But what content defines ``not
knowing''? Is there empirical evidence about the unknown? The unknown
can only be an abstract concept, expressed through probability theory or
a logical equivalent. Is there empirical evidence that probabilities are
reducible, that there is a simpler logical equivalent? No, quite the
opposite.

The only trade-off seems to be between costs (time complexity, etc.) and
understanding/control. A prior measure which we understand and/or
control may mean much more costs than an emergent (thus, inconsistent
and uncontrollable) prior measure which just minimizes some distance.
But this trade-off is not new, and it is already present in all
industries which deal with some safety risk (which is essentially all
industries). Distances are efficient for proof of concepts (pilot
projects), when the goal is to show that we are a short distance from
where we want to be. But safety (as most features) is not being at a
short distance from being safe\footnote{See for instance
  \href{https://edition.cnn.com/2023/04/29/us/ai-scam-calls-kidnapping-cec/index.html}{https://edition.cnn.com/2023/04/29/us/ai-scam-calls-kidnapping-cec}}.
``We were at a short distance from avoiding nuclear annihilation'' is
completely different from ``we avoided nuclear annihilation''. To avoid
nuclear annihilation we need (probability) measures, not only distances.

\hypertarget{status-of-the-machine-checked-formalization}{%
\section{10. Status of the machine-checked formalization (Lean
4)}\label{status-of-the-machine-checked-formalization}}

The mathematical content of this article has been formalized in Lean~4
(with Mathlib) in the \texttt{PnpProof} library of this repository; the
formalization plan is \texttt{PNP\_IMPLEMENTATION\_PLAN.md} and the
per-lemma map is \texttt{PnpProof/IMPLEMENTED.md}. The development
compiles \emph{without any \texttt{sorry} and without any axioms} beyond
Lean's standard three (\texttt{propext}, \texttt{Classical.choice},
\texttt{Quot.sound}); the build was last verified on 2026-06-17.

The following machine-checked theorems cover the claims of this article:

\begin{itemize}
\item
  \texttt{selection\_not\_history} (Sections 1, 3): selection by
  conditioning is well-defined (via disintegration into regular
  conditional probabilities), every target output is a null event, and no
  countable history of events realizes the selection --- the fact in the
  title.
\item
  \texttt{almost\_all\_not\_computable} (Sections 2, 3, 6): with prior
  probability one, the selected function is computed by \emph{no}
  deterministic machine (at any time cost --- an information-theoretic
  statement stronger than ``not polynomial-time''). The supporting chain
  is fully formalized: computable codes are countable, countable sets are
  null for the (proved atomless) prior.
\item
  The counting results of Section 2: Boolean circuits are defined from
  scratch and Shannon's bound is proved --- almost all Boolean functions
  on \(n\) bits (all but a fraction \(2^{-2^{n}/4}\)) need more than
  \(2^{n}/(4n)\) gates, for \(n \geq 8\).
\item
  The measure-theoretic facts of Section 5: a finite measure has
  countably many atoms, continuous and atomic parts are mutually
  singular, a jump separates a mixed cumulative distribution from a
  continuous one on a rational interval, and conditioning a mixed measure
  on its diffuse part yields an atomless probability measure.
\item
  The uniform measure on the infinite-dimensional sphere (Sections 2, 3):
  the rescaled-Gegenbauer-to-Hermite limit of L\'opez and Temme (proved
  by recurrence induction), the weight and normalization-constant limits,
  the uniform measure on the \(k\)-sphere realized as a normalized
  Gaussian with full rotation invariance, and the Poincar\'e--Borel
  theorem (sphere marginals become Gaussian; almost-sure concentration on
  the sphere via the strong law of large numbers).
\item
  \texttt{model\_P\_ne\_NP} (Section 6, option 1): within the model of
  Sections 3--6, verification of a candidate output is a computable
  (cheap) operation, while no deterministic machine computes the selected
  function, prior-almost-surely.
\item
  \texttt{model\_P\_ne\_NP\_circuit}: the same separation with the
  verification side strengthened to an explicit Boolean comparator
  circuit of at most \(50k + 50\) gates on the three \(k\)-bit dyadic
  inputs (a big-endian ripple comparator, built gate by gate with a
  machine-checked correctness proof).
\item
  \texttt{model\_vs\_clay\_disjointness}: for \emph{any} collection of
  languages all of whose members are computable --- as is every faithful
  formalization of the Clay classes \(P\) and \(NP\), since
  \(NP \subseteq EXPTIME\) --- prior-almost-surely no member of the
  collection decides the selected function's dyadic data.
\item
  \texttt{interpPi02\_geom\_iff}: every \(\Pi^0_2\) arithmetical sentence
  \(\forall x\,\exists y\,p(x,y)\) is encoded as a single completed vector
  \(\Psi_p = \sum_x [\mathrm{OK}_p(x)]/(x+1)\,e_x\) in the separable Hilbert
  substrate of a model, and its truth is the metric identity
  \(\lVert \Psi_p - \Psi_{\mathrm{true}}\rVert = 0\); this identity is proved
  \emph{equal} to the sentence. The coordinate \(\mathrm{OK}_p\) is
  noncomputable (it decides the inner \(\Sigma^0_1\) formula), so this is an
  \emph{expressibility} device --- it stores the answer, it does not compute it
  --- and it yields no implication to the Clay statement.
\item
  \texttt{npc\_not\_inP}: the NP-completeness closure theorem --- a single
  language in \(NP \setminus P\) forces every NP-complete problem (e.g.\ SAT)
  out of \(P\) --- proved over an abstract complexity model whose one
  load-bearing property is that \(P\) is closed under polynomial-time
  reductions. (A faithful concrete runtime model is left abstract because
  Mathlib's fuel-indexed evaluator bounds a recursion depth, not a step count;
  whether the present construction supplies the \(NP \setminus P\) witness is
  exactly the open modelling question above.)
\end{itemize}

To keep the formal development honest, the following are
\emph{deliberately not} formalized. The claim that the model separation
implies the standard (Clay Millennium) \(P \neq NP\) statement is a
modelling argument about what counts as the proper computer model; as
stated in the Abstract, it still needs to be reviewed by experts in
Complexity Theory, and accordingly it appears in the formalization only
as prose --- the theorem is named \texttt{model\_P\_ne\_NP}, never
\texttt{P\_ne\_NP}, and no implication between the two is asserted in
Lean. The \emph{relationship} between the two statements is itself
machine-checked: \texttt{model\_vs\_clay\_disjointness} proves that the
model's hard object lies, almost surely, outside the arena of computable
languages in which the Clay problem is played --- which is why neither
statement implies the other. Section 7 (the realistic version) is a
sketch: a purely analytic
approximation statement would be false (polynomials are dense in
\(C([0,1])\) and \(L^{2}([0,1])\)), so the intended
\emph{computational} approximation statement requires a precise
formulation before it can be formalized; no statement has been committed.
Section 8 is, as noted there, not mathematically provable, and Section 9
is discussion.

All mathematical items queued in the plan are now proved; the only
remaining item is minor build-lint cleanup.